%% Preamble for the FastVCL CVPR draft.
%% The document class and cvpr.sty are loaded in main.tex; this file only adds
%% packages, math macros and the visible TODO marker.

% ---------------------------------------------------------------------------
% Layout / tables / algorithms
% ---------------------------------------------------------------------------
\usepackage{array}
\usepackage{multirow}
\usepackage{tabularx}
\usepackage{algorithm}
\usepackage{algpseudocode}
\usepackage{enumitem}
\usepackage{amsthm}

% Algorithm floats inside the two-column text
\renewcommand{\algorithmicrequire}{\textbf{Input:}}
\renewcommand{\algorithmicensure}{\textbf{Output:}}
\floatname{algorithm}{Algorithm}

% Slightly tighter lists
\setlist[itemize]{leftmargin=1.2em,topsep=2pt,itemsep=1pt,parsep=0pt}
\setlist[enumerate]{leftmargin=1.5em,topsep=2pt,itemsep=1pt,parsep=0pt}

% ---------------------------------------------------------------------------
% Visible placeholders.  DO NOT redefine these to be empty: an empty TODO would
% make the draft look finished while the evidence is still outstanding.
% ---------------------------------------------------------------------------
\newcommand{\red}[1]{{\color{red}#1}}
\newcommand{\todo}[1]{{\color{red}#1}}
\newcommand{\TODO}[1]{\textbf{\color{red}[TODO: #1]}}
\newcommand{\PENDING}{\TODO{PENDING}}
\newcommand{\pending}{--}

% ---------------------------------------------------------------------------
% Math notation
% ---------------------------------------------------------------------------
\newcommand{\fA}[1]{\bar{f}^{A}_{#1}}
\newcommand{\fB}[1]{\bar{f}^{B}_{#1}}
\newcommand{\Sim}{S}
\newcommand{\Tube}{\mathcal{T}}
\newcommand{\Anch}{\mathcal{A}}
\newcommand{\Grid}{\mathcal{G}}
\newcommand{\Pred}{\mathcal{P}}
\newcommand{\Step}{\mathcal{S}}
\newcommand{\Track}{\mathcal{I}}
\newcommand{\TopK}{\operatorname{TopK}}
\newcommand{\Dedup}{\operatorname{Dedup}}
\newcommand{\AnchorCap}{\operatorname{Cap}}
\newcommand{\minIoU}{\operatorname{IoU}_{\min}}
\newcommand{\IoU}{\operatorname{IoU}}
\newcommand{\Tloc}{T_{\mathrm{localize}}}
\newcommand{\Taug}{T_{\mathrm{aug}}}
\DeclareMathOperator*{\argmax}{arg\,max}
\DeclareMathOperator*{\argmin}{arg\,min}

% Shorthand for the two video labels (never use them as a batch size)
\newcommand{\vidA}{\mathrm{A}}
\newcommand{\vidB}{\mathrm{B}}

% ---------------------------------------------------------------------------
% Small spacing helpers
% ---------------------------------------------------------------------------
\newcommand{\tabnote}[1]{\par\vspace{-2pt}{\footnotesize #1}}
\setlength{\abovedisplayskip}{4pt}
\setlength{\belowdisplayskip}{4pt}
\setlength{\abovedisplayshortskip}{2pt}
\setlength{\belowdisplayshortskip}{2pt}
